\section{Introduction}
\label{sec:intro}

\IEEEPARstart{A}{hospital} refers a patient to a home-care provider. The provider accepts and starts care on time, but its acceptance reply is lost, its start notice is lost, and when the hospital reads the provider's record the read is served from a replica that has not caught up. Every message the protocol exchanged was well formed, every transition the provider took was legal, and the obligation the provider owed was discharged. Yet a hospital that judges by its protocol concludes that the provider never accepted. Under a choreography, a state machine and a commitment protocol alike, the hospital's judgment is wrong; under one hybrid design it is right about care and declines to judge the acceptance (\cref{sec:example}).

The example separates three questions that the literature on coordination protocols usually answers together. Did the dependency the protocol exists to manage actually hold? Do the participant's observations determine whether it held? Does the participant's judgment say so correctly? Communication protocols are verified by asking whether interactions conform to a global specification~\cite{honda2016mpst,scalas2019less}. State-machine protocols are verified by asking whether permitted transitions preserve invariants~\cite{harel1987statecharts,lynch1989ioa}. Commitment protocols are verified by asking whether social obligations are created, managed and discharged~\cite{singh1999ontology,yolum2002flexible}. Each question is answered relative to a paradigm's own semantics, and none of them is the question a coordinating organization asks: whether the dependency between its activities was managed, and whether it can establish that it was~\cite{malone1994interdisciplinary}.

\subsection{Thesis}

This paper argues that protocol adequacy is a relation, not a property, and that it factors. A protocol is adequate \emph{for} a dependency, \emph{under} an environment of admissible faults, \emph{with respect to} an evidence model. Adequacy then decomposes into three parts that can be measured separately: \emph{conduct}, whether the protocol's participants act so that the dependency holds; \emph{observability}, whether the observations the protocol gives a participant determine whether it holds; and \emph{interpretation}, whether the participant's judgment over those observations is sound. The second part has an exact characterization. For any protocol and environment, the largest fraction of judgments a sound judge can make definitively equals the probability mass of observations that determine the truth (\cref{thm:ceiling}). Epistemic adequacy is therefore a property of what a protocol makes observable; a judge can at best attain it and at worst betray it.

\subsection{Contributions}

\begin{enumerate}
\item \textbf{An adequacy calculus} (\cref{sec:systems,sec:calculus}). Coordination systems, dependencies and protocols are defined as formal objects; six axioms state the theory's commitments; operational and epistemic adequacy are separated. Four results are proved: indistinguishable observations preclude sound complete judgment (\cref{thm:p1}); determinable mass is the exact ceiling on sound coverage and is attained by a canonical judge (\cref{thm:ceiling}); a composition is compliance-decidable for an observer only if every relay on the observer's path distinguishes the truth-relevant states (\cref{thm:p2}); and refinement preserves adequacy only under four separately necessary conditions (\cref{thm:p3}).
\item \textbf{An executable-semantics method} (\cref{sec:method}). Each scenario family is a closed PLT Redex language~\cite{felleisen2009redex}; a world model alone holds oracle truth; arms receive only their deliveries, enforced by a check on every arm's import closure; a judge decides determinacy by exhaustive search over each environment's support, so every metric is an exact expectation, and seeded sampling with common random numbers is checked against it.
\item \textbf{DMBench results} (\cref{sec:dmbench,sec:results}). Eight scenario families, five arms (choreography, state machine, commitments, a hybrid with an evidentiary discipline, and that hybrid's ablation) and thirty-six fault profiles. No single paradigm is safe across the families; the hybrid makes no false claim under any profile; its ablation is unsafe in every family; false completion follows the evidence model, not the paradigm; and operational and epistemic adequacy separate in four families.
\item \textbf{A differential protocol for independent authorship} (\cref{sec:independent}). Arms rewritten by independent authors from written specifications alone are compared with the reference arms on every reachable observation sequence. Adjudication against the specification reduced 310 disagreements to 33 documented interpretive choices and found twelve errors in the reference arms; the hybrid's safety and its ablation's unsafety hold for the independent arms before any adjudication.
\end{enumerate}

\subsection{What the paper does not claim}

The paper does not claim that coordination, commitments, session types or protocol composition are new; it builds on each. It does not claim that one paradigm wins: the strongest result of a comparative benchmark is a map of which dependency and environment combinations favour which mechanism~\cite{malone1994interdisciplinary}. It does not claim statistical generalization to deployed systems: the models are finite, discrete and exhaustively enumerated, which is what makes their metrics exact and what limits their reach (\cref{sec:limitations}). The ontology and the protocol-selection method proposed with the framework are stated as research directions, not results (\cref{sec:agenda}).

\subsection{Organization}

\Cref{sec:related} places the work. \Cref{sec:systems} defines coordination systems, dependencies and protocols. \Cref{sec:calculus} develops the adequacy calculus. \Cref{sec:method} describes the executable-semantics method and \cref{sec:dmbench} the benchmark. \Cref{sec:independent} presents the differential protocol. \Cref{sec:results} reports results against the framework's hypotheses, \cref{sec:example} carries one episode through every layer, and \cref{sec:implications} states the design rules the results support. \Cref{sec:limitations,sec:agenda} state limitations and open problems. Appendices give full proofs, family specifications, the adjudication record and the complete exact results.
